\documentclass[11pt,a4paper]{article}
\usepackage[utf8]{inputenc}
\usepackage[vietnamese]{babel}
\usepackage{amsmath,amssymb}
\usepackage{geometry}
\geometry{a4paper, total={180mm,260mm}, left=15mm, right=15mm, top=18mm, bottom=18mm}
\usepackage{tikz}
\usetikzlibrary{shapes.geometric, arrows.meta, positioning, calc}
\usepackage{tcolorbox}
\tcbuselibrary{skins,breakable}
\usepackage{xcolor}
\usepackage{multicol}
\usepackage{array}
\usepackage{fontawesome}
\usepackage{fancyhdr}

\pagestyle{fancy}
\fancyhf{}
\renewcommand{\headrulewidth}{0.5pt}
\renewcommand{\footrulewidth}{0.5pt}
\fancyhead[L]{\small\bfseries\color{teal}\faBook\ SỔ TAY KIẾN THỨC TOÁN 3}
\fancyhead[R]{\small\bfseries\color{teal}Lớp Toán Thầy Lĩnh -- loptoanthaylinh.vn}
\fancyfoot[C]{\small\color{gray}-- Sổ tay kiến thức: Bài 35. Luyện tập chung đo lường (mm, g, ml, $^\circ\text{C}$) --}

\definecolor{primary}{HTML}{005F73}
\definecolor{secondary}{HTML}{0A9396}
\definecolor{accent}{HTML}{EE9B00}
\definecolor{danger}{HTML}{AE2012}
\definecolor{boxbg}{HTML}{F4F6F8}

\begin{document}

% --- BANNER TIÊU ĐỀ ---
\begin{tcolorbox}[
  enhanced,
  colback=white,
  colframe=primary,
  arc=3mm,
  boxrule=1.2pt,
  left=4mm, right=4mm, top=2.5mm, bottom=2.5mm,
  drop shadow=black!10!white
]
  \centering
  {\Large\bfseries\color{primary}\faBookmark\ BÀI 35: LUYỆN TẬP CHUNG ĐO LƯỜNG}\\[1.5mm]
  {\small\itshape\color{black!75}Tổng hợp kĩ năng đổi đơn vị, tính toán kèm tên đơn vị và giải toán đố thực tế trên 1 trang A4}
\end{tcolorbox}

\vspace{-2mm}


\noindent
\begin{minipage}[t]{0.49\textwidth}\vspace{0pt}
\begin{tcolorbox}[
  enhanced,
  colback=blue!4!white,
  colframe=blue!75!black,
  arc=2mm, boxrule=0.8pt,
  left=2mm, right=2mm, top=1.8mm, bottom=1.8mm,
  title={\footnotesize\bfseries\color{white}\faKey\ 1. Ma trận 4 bảng đổi đơn vị cốt lõi}
]
{\footnotesize
\begin{itemize}\setlength{\itemsep}{1.5pt}\setlength{\leftskip}{-2.5mm}
  \item \textbf{Độ dài:} $1\text{ cm} = 10\text{ mm}$; $1\text{ m} = 1000\text{ mm}$; $1\text{ km} = 1000\text{ m}$.
  \item \textbf{Khối lượng:} $1\text{ kg} = 1000\text{ g} \iff 1000\text{ g} = 1\text{ kg}$.
  \item \textbf{Dung tích:} $1\ l = 1000\text{ ml} \iff 1000\text{ ml} = 1\ l$.
  \item \textbf{Nhiệt độ:} Đo bằng $^\circ\text{C}$; Thân nhiệt $37^\circ\text{C}$; Nước sôi $100^\circ\text{C}$.
  \item Các tiền tố "kilo" $\implies 1000$; "mili" $\implies \dfrac{1}{1000}$.
\end{itemize}
}
\end{tcolorbox}
\end{minipage}
\hfill
\begin{minipage}[t]{0.49\textwidth}\vspace{0pt}
\begin{tcolorbox}[
  enhanced,
  colback=teal!4!white,
  colframe=teal!75!black,
  arc=2mm, boxrule=0.8pt,
  left=2mm, right=2mm, top=1.8mm, bottom=1.8mm,
  title={\footnotesize\bfseries\color{white}\faCheckSquare\ 2. Kĩ năng thực hiện phép tính}
]
{\footnotesize
\begin{itemize}\setlength{\itemsep}{1.5pt}\setlength{\leftskip}{-2.5mm}
  \item \textbf{Cộng / Trừ:} Bắt buộc đổi về \textbf{cùng một đơn vị} rồi thực hiện phép tính bình thường, viết kèm đơn vị vào kết quả.
  \item \textbf{Nhân với một số:}
  $$A\ (\text{đơn vị}) \times k = (A \times k)\ (\text{đơn vị})$$
  \item \textbf{Chia cho một số:}
  $$A\ (\text{đơn vị}) : k = (A : k)\ (\text{đơn vị})$$
\end{itemize}
}
\end{tcolorbox}
\end{minipage}

\vspace{-2.5mm}
% --- KHỐI 3: SƠ ĐỒ TƯ DUY TIKZ ---
\noindent
\begin{tcolorbox}[
  enhanced,
  colback=orange!4!white,
  colframe=orange!80!black,
  arc=2mm, boxrule=0.8pt,
  left=2mm, right=2mm, top=1.5mm, bottom=1.5mm,
  title={\footnotesize\bfseries\color{white}\faSitemap\ SƠ ĐỒ TƯ DUY: HỆ ĐO LƯỜNG TOÀN DIỆN LỚP 3}
]
\centering
\begin{tikzpicture}[
  box/.style={draw, rounded corners=2pt, align=center, font=\scriptsize\bfseries, inner sep=2.5pt, minimum height=0.75cm},
  arr/.style={-{Stealth[length=2mm]}, thick}
]
  \node[box, fill=purple!15, draw=purple!80!black, text width=2.6cm] (p1) at (0, 0) {ĐỘ DÀI\\$1\text{ cm} = 10\text{ mm}$\\$1\text{ m} = 1000\text{ mm}$};
  \node[box, fill=teal!15, draw=teal!80!black, text width=2.6cm] (p2) at (3.5, 0) {KHỐI LƯỢNG\\$1\text{ kg} = 1000\text{ g}$\\$500\text{ g} = \frac{1}{2}\text{ kg}$};
  \node[box, fill=blue!12, draw=blue!80!black, text width=2.6cm] (p3) at (7.0, 0) {DUNG TÍCH\\$1\ l = 1000\text{ ml}$\\$250\text{ ml} = \frac{1}{4}\ l$};
  \node[box, fill=orange!15, draw=orange!80!black, text width=2.6cm] (p4) at (10.5, 0) {NHIỆT ĐỘ\\$^\circ\text{C}$\\Nhiệt kế đo};

  \draw[arr, purple!80!black] (p1) -- (p2);
  \draw[arr, teal!80!black] (p2) -- (p3);
  \draw[arr, blue!80!black] (p3) -- (p4);
\end{tikzpicture}
\end{tcolorbox}

\vspace{-2.5mm}
% --- KHỐI 4: VÍ DỤ & MẸO TRÁNH BẪY ---
\noindent
\begin{minipage}[t]{0.49\textwidth}\vspace{0pt}
\begin{tcolorbox}[
  enhanced,
  colback=green!4!white,
  colframe=green!60!black,
  arc=2mm, boxrule=0.8pt,
  left=2mm, right=2mm, top=1.8mm, bottom=1.8mm,
  title={\footnotesize\bfseries\color{white}\faCheckCircle\ Ví dụ thực hành chuẩn mực}
]
{\footnotesize
\textbf{1. Thực hiện phép tính:}\\
$150\text{ g} + 350\text{ g} = \mathbf{500\text{ g}} = \mathbf{\dfrac{1}{2}\text{ kg}}$.\\
$1\ l - 300\text{ ml} = 1000\text{ ml} - 300\text{ ml} = \mathbf{700\text{ ml}}$.\\
$24\text{ mm} : 4 = \mathbf{6\text{ mm}}$.\\
\textbf{2. Bài toán 2 bước có đơn vị đo:}\\
Can thứ nhất có $15\ l$ dầu, can thứ hai có số lít dầu gấp đôi can thứ nhất. Cả hai can có:\\
Can hai: $15 \times 2 = 30\ (l)$. Cả hai can: $15 + 30 = \mathbf{45}\ (l)$.
}
\end{tcolorbox}
\end{minipage}
\hfill
\begin{minipage}[t]{0.49\textwidth}\vspace{0pt}
\begin{tcolorbox}[
  enhanced,
  colback=red!4!white,
  colframe=red!70!black,
  arc=2mm, boxrule=0.8pt,
  left=2mm, right=2mm, top=1.8mm, bottom=1.8mm,
  title={\footnotesize\bfseries\color{white}\faExclamationTriangle\ Mẹo tránh bẫy trắc nghiệm}
]
{\footnotesize
\begin{itemize}\setlength{\itemsep}{1.5pt}\setlength{\leftskip}{-2.5mm}
  \item \textbf{Bẫy quên ghi tên đơn vị đo:}\\
  Khi làm bài tự luận hoặc điền đáp án, nếu kết quả có đơn vị mà chỉ ghi con số (ví dụ $45$ thay vì $45\ l$) sẽ bị mất điểm hoàn toàn!
  \item \textbf{Bẫy nhân chia cả đơn vị:}\\
  $10\text{ cm} \times 2 = 20\text{ cm}$ (đúng), KHÔNG viết thành $20\text{ cm}^2$ (đây là diện tích)!
\end{itemize}
}
\end{tcolorbox}
\end{minipage}


\end{document}
